\section*{第 \ref{ch:b1:predominant-reaction} 章——酸碱平衡与主反应法}

\begin{solution}{exo:b1:predominant-reaction:1}
共轭碱：\ce{HSO4-}、\ce{CO3^{2-}}、\ce{NH3}、\ce{OH-}。共轭
酸：\ce{NH4+}、\ce{H2PO4-}、\ce{H2O}、\ce{H3O+}。\omterm{def:b1:predominant-reaction:ampholyte}{两性物质}：\ce{HCO3-}、
\ce{H2O}、\ce{HPO4^{2-}}（还有 \ce{HSO4-}，即硫酸的共轭碱，
以及其第二对共轭酸碱中的酸）。
\end{solution}

\begin{solution}{exo:b1:predominant-reaction:2}
\ce{HF} ($K_a = \num{6.9e-4}$) $>$ \ce{CH3COOH} (\num{1.7e-5}) $>$ \ce{HClO}
(\num{2.8e-8}) $>$ \ce{NH4+} (\num{5.6e-10}).
\end{solution}

\begin{solution}{exo:b1:predominant-reaction:3}
\omterm{def:b1:predominant-reaction:strong-acid}{强酸}：$\mathrm{pH} = 2.00$。\omterm{def:b1:predominant-reaction:strong-acid}{强碱}：$[\ce{OH-}] = 10^{-2}$，
$\mathrm{pH} = 14.00 - 2.00 = 12.00$.
\end{solution}

\begin{solution}{exo:b1:predominant-reaction:4}
\ce{CH3COOH} 在 pH 低于 4.76 时占优，\ce{CH3COO-} 在高于此值时占优。pH 为 3 时，酸型
占优（98\,\%）；在血液（pH 7.4）中则为乙酸根离子。pH 为 3.76 时，
酸型分数为 $1/(1 + 10^{-1}) = 0.91$。
\end{solution}

\begin{solution}{exo:b1:predominant-reaction:5}
主反应为 \ce{NH3 + H2O <=> NH4+ + OH-}，$K = 10^{-(14.00 -
9.25)} = 10^{-4.75}$。若只有少量 \ce{NH3} 反应，则 $[\ce{OH-}] = \sqrt{K C} =
\qty{1.3e-3}{mol/L}$，$\mathrm{pH} = 11.13$。验证：$[\ce{NH4+}] = 1.3\,\%$
为 $C$ 的比例，且 $\mathrm{pH} > \mathrm{p}K_a + 1$，故成立。
\end{solution}

\begin{solution}{exo:b1:predominant-reaction:6}
$\mathrm{pH} = \frac12(4.76 + 2.00) = 3.38$（精确值为 3.39）；
$\alpha = 10^{-3.38}/0.010 = 0.042$：解离 4\,\%，低于 10\,\%，且
$\mathrm{pH} < \mathrm{p}K_a - 1$，故成立。
\end{solution}

\begin{solution}{exo:b1:predominant-reaction:7}
\ce{HF + OH- -> F- + H2O}，$K = 10^{14.00 - 3.16} = 10^{10.84}$：完全反应。
\omterm{def:b1:predominant-reaction:predominant-reaction}{等效溶液}中 \ce{HF} 与 \ce{F-} 各为 \qty{0.050}{mol/L}，构成
\omterm{def:b1:predominant-reaction:buffer}{缓冲溶液}：$\mathrm{pH} = \mathrm{p}K_a = 3.16$。
\end{solution}

\begin{solution}{exo:b1:predominant-reaction:8}
$\mathrm{pH} = 4.76$。盐酸将 \qty{0.010}{mol}
乙酸根转化为乙酸：$\mathrm{pH} = 4.76 + \log(0.040/0.060) = 4.58$，变化量为
0.18。在纯水中，pH 将由 7.00 降至 2.00。
\end{solution}

\begin{solution}{exo:b1:predominant-reaction:9}
在水中两种酸都完全反应生成 \ce{H3O+}：溶剂将
它们拉平。乙酸的碱性远弱于水，因此在乙酸中两种酸
只部分反应，且程度不同，
从而可以比较其强弱。水中能够存在的最\omterm{def:b1:predominant-reaction:strong-acid}{强酸}是 \ce{H3O+}。
\end{solution}

\begin{solution}{exo:b1:predominant-reaction:10}
\ce{CO3^{2-} + H2O <=> HCO3- + OH-}，$K = 10^{-(14.00 - 10.33)} = 10^{-3.67}$；
$[\ce{OH-}] = \sqrt{10^{-3.67} \times 0.10} = \qty{4.6e-3}{mol/L}$（4.6\,\%
占 $C$ 的），$\mathrm{pH} = 11.67$。第二步碱性反应，\ce{HCO3- + H2O <=>
CO2 + H2O + OH-}，$K = 10^{-7.63}$：它给出 $[\ce{CO2}] \approx
10^{-7.6}$，可忽略。
\end{solution}

\begin{solution}{exo:b1:predominant-reaction:11}
$[\mathrm{A}] = [\ce{H3O+}] = \alpha C$ 且 $[\mathrm{HA}] = (1 - \alpha)C$，
故 $K_a = \alpha^2C/(1 - \alpha)$。取 $C = 10^{-3}$，得 $\alpha^2 +
0.0174\,\alpha - 0.0174 = 0$，$\alpha = 0.12$。该酸解离 12\,\%，
因此少量解离的假设（$\alpha <
10\,\%$）不成立，须求解二次方程。
\end{solution}

\begin{solution}{exo:b1:predominant-reaction:12}
盐酸使 $[\ce{H3O+}] = \qty{0.010}{mol/L}$，pH 为 2.00。于是
\[
  [\ce{CH3COO-}] = \frac{K_a[\ce{CH3COOH}]}{[\ce{H3O+}]} = \frac{1.74 \times 10^{-5}
  \times 0.10}{0.010} = \qty{1.7e-4}{mol/L},
\]
与 0.010 相比可以忽略，
来自\omterm{def:b1:predominant-reaction:strong-acid}{强酸}的 \ce{H3O+} 使\omterm{def:b1:predominant-reaction:strong-acid}{弱酸}的平衡
逆向移动。
\end{solution}

\begin{solution}{pb:b1:predominant-reaction:1}
\textbf{1.} \ce{H3PO4}/\ce{H2PO4-}、\ce{H2PO4-}/\ce{HPO4^{2-}}、
\ce{HPO4^{2-}}/\ce{PO4^{3-}}；例如 \ce{H3PO4 <=> H2PO4- + H+}。
\textbf{2.} 分界点位于 2.15、7.21 和 12.34。
\textbf{3.} pH 为 1 时为 \ce{H3PO4}；5 时为 \ce{H2PO4-}；9 时为 \ce{HPO4^{2-}}；13 时为
\ce{PO4^{3-}}。
\textbf{4.} \ce{H2PO4-} 与 \ce{HPO4^{2-}}。
\textbf{5.} 每失去一个质子，就留下一个负电荷更高的离子，
它更牢固地结合下一个质子。
\textbf{6.} \ce{H3PO4 + H2O <=> H2PO4- + H3O+}；公式
$\frac12(\mathrm{p}K_a + \mathrm{p}C) = 1.58$ 不满足 $\mathrm{pH} <
\mathrm{p}K_a - 1 = 1.15$：约五分之一的酸发生反应。
\textbf{7.} $x^2 + K_{a1}x - K_{a1}C = 0$，其中 $K_{a1} = 10^{-2.15}$，得 $x =
\qty{0.0233}{mol/L}$，$\mathrm{pH} = 1.63$。
\textbf{8.} \ce{2H2PO4- <=> H3PO4 + HPO4^{2-}}，$K = 10^{2.15 - 7.21} =
10^{-5.06}$.
\textbf{9.} $\mathrm{pH} = \frac12(2.15 + 7.21) = 4.68$。
\textbf{10.} $\mathrm{pH} = \frac12(7.21 + 12.34) = 9.78$。
\textbf{11.} 无：1.63、4.68 和 9.78 都不能单独构成\omterm{def:b1:predominant-reaction:buffer}{缓冲溶液}，需要
其中两种的混合物。
\textbf{12.} \ce{H3PO4 + OH- -> H2PO4- + H2O}，$K = 10^{14.00 - 2.15} =
10^{11.85}$：\omterm{def:b1:extent-q-and-k:quantitative}{定量反应}。
\textbf{13.} \qty{0.100}{mol/L} 的 \ce{H2PO4-}：pH 为 4.68。
\textbf{14.} 再加入 \qty{0.050}{mol} 氢氧根，可使其中一半转化：
\qty{0.050}{mol/L} 的 \ce{H2PO4-} 与 \ce{HPO4^{2-}}，pH 为 7.21。
\textbf{15.} \qty{0.100}{mol/L} 的 \ce{HPO4^{2-}}：pH 为 9.78。
\textbf{16.} pH 自 1.63 升高，在 $n = 0.100$ 附近跃升（至 4.68），
在 $n = 0.150$ 处缓慢经过 7.21，在 $n = 0.200$ 附近跃升（至
9.78），在 $n = 0.250$ 附近经过 12.34，在 $n = 0.300$ 时达到约 12.6（\qty{0.100}{mol/L} 的 \ce{PO4^3-}：$x^2/(0.100 - x) = 10^{-1.66}$，$x = 0.037$，pH 12.57）。
\textbf{17.} 它含等量的酸及其共轭碱，即
$\mathrm{p}K_a$ 为 7.21 的\omterm{def:b1:predominant-reaction:bronsted}{酸碱共轭对}（见\cref{prop:b1:predominant-reaction:buffer}）。
\textbf{18.} $10^{7.20 - 7.21} = 0.977$.
\textbf{19.} $n(\ce{H2PO4-}) = 0.100 - x$, $n(\ce{HPO4^{2-}}) = x$.
\textbf{20.} $x/(0.100 - x) = 0.977$, $x = \qty{0.0494}{mol}$.
\textbf{21.} 酸使 \qty{0.0010}{mol} 的 \ce{HPO4^{2-}} 质子化：
$\mathrm{pH} = 7.21 + \log(0.0484/0.0516) = 7.18$，pH 变化量为 $-0.02$。
在纯水中，pH 由 7.00 变为 3.00。
\textbf{22.} 否：pH 取决于物质的量之比，而稀释
不改变这一比值（只要\omterm{def:b1:extent-q-and-k:activity}{活度}仍接近
浓度）。
\textbf{23.} 所需氢氧根 $0.100 + 0.0494 = \qty{0.149}{mol}$：\textbf{\qty{149}{mL}}
取 \qty{1.00}{mol/L} 溶液，配制成 \qty{1.00}{L}。
\end{solution}
